% Generated draft: requires mathematical and source review.
\documentclass[11pt]{article}
\usepackage[margin=25mm]{geometry}
\usepackage{fontspec}
\setmainfont{Latin Modern Roman}
\usepackage{amsmath,amssymb,mathtools,unicode-math}
\setmathfont{Latin Modern Math}
\usepackage{enumitem,needspace,xurl,hyperref}
\setlength{\parindent}{0pt}
\setlength{\parskip}{5pt}
\setlength{\emergencystretch}{3em}
\newcommand{\N}{\mathbb{N}}
\newcommand{\Z}{\mathbb{Z}}
\newcommand{\Q}{\mathbb{Q}}
\newcommand{\R}{\mathbb{R}}
\newcommand{\C}{\mathbb{C}}
\newcommand{\F}{\mathbb{F}}
\newcommand{\A}{\mathbb{A}}
\newcommand{\PP}{\mathbb{P}}
\newcommand{\E}{\mathbb{E}}
\newcommand{\eps}{\varepsilon}
\newcommand{\dd}{\,\mathrm{d}}
\begin{document}
\section*{UnsolvedMath ID: 30001737}
\textbf{Source problem code:} catalogue entry
\subsection{Definitions and mathematical statement}
The supplied record asks for a converse to a unitary distinction criterion. The following formulation uses the customary ambient interpretation, which is not fully specified in the record. Let \(p,q\) be signature parameters, put \(n=p+q\), and set
\[
G=\operatorname{GL}_n(\C),\qquad J_{p,q}=\operatorname{diag}(I_p,-I_q),\qquad H=U(p,q)=\{g\in G:\overline g^{\,t}J_{p,q}g=J_{p,q}\}.
\]
Here signature parameters are conventionally nonnegative integers; the supplied statement does not specify whether zero entries or rank zero are included.

Let \(\boldsymbol\chi=(\chi_1,\ldots,\chi_n)\) be characters of \(\C^\times\) forming Langlands standard data, and let \(\pi\) be the corresponding irreducible Langlands quotient. Conventionally, this means the unique irreducible quotient of the normalized principal series induced from these characters, placed in Langlands order, from the upper triangular Borel subgroup of \(G\). The precise character hypotheses and representation category are not supplied.

Write \(\pi^c(g)=\pi(\overline g)\). Interpret conjugation invariance as \(\pi\simeq\pi^c\), and interpret \(H\)-distinction as
\[
\operatorname{Hom}_{H}(\pi,\mathbf 1)\ne\{0\},
\]
where \(\mathbf 1\) is the trivial representation and the continuity requirements for the invariant functional must be fixed by the intended representation category.

For a character \(\chi\), write \(\chi^c(z)=\chi(\overline z)\). Under the interpretation that conjugation acts on the Langlands character multiset, let \(m(\chi)\) denote its multiplicity and let \(S\) be its support. When this multiset is conjugation invariant, two possible readings of the orbit count are
\[
 k_{\mathrm{distinct}}=\#\{\{\chi,\chi^c\}:\chi\in S,\ \chi\ne\chi^c\},
 \qquad
 k_{\mathrm{mult}}=\frac12\sum_{\substack{\chi\in S\\\chi\ne\chi^c}}m(\chi).
\]
The first counts distinct two-element character orbits; the second counts conjugate pairs with multiplicity. The supplied record does not identify the underlying orbit set or settle this multiplicity convention. Denote its intended count by \(k_2(\boldsymbol\chi)\).

\textbf{Requested converse, subject to the unresolved conventions above.} For every permitted signature \((p,q)\) and every Langlands quotient \(\pi\) of the indicated characters, does
\[
\left(\pi\simeq\pi^c\ \text{ and }\ k_2(\boldsymbol\chi)\leq\min(p,q)\right)
\quad\Longrightarrow\quad
\operatorname{Hom}_{U(p,q)}(\pi,\mathbf 1)\ne\{0\}
\]
hold without assuming that \(\pi\) is generic? Here generic conventionally means admitting a nonzero Whittaker functional for a nondegenerate character of the maximal upper unipotent subgroup. The supplied source asserts the reverse implication as a necessary condition and presents the displayed converse as the conjecture. A fully unambiguous statement requires confirmation of the ambient conventions, admissible signatures, and orbit-count definition; the two counts displayed above are alternatives, not two separately supplied conjectures.
\subsection{Short English statement}
The question asks whether conjugation invariance and a bound of \(\min(p,q)\) on the number of conjugate pairs in the Langlands character data guarantee a nonzero \(U(p,q)\)-invariant functional for every such Langlands quotient, including nongeneric ones. The precise orbit-count convention and representation setting need clarification. The supplied partially solved label is dated curation, not an independently verified conclusion.
\subsection{Sources}
\begin{itemize}
\item Local registry: \texttt{data/unsolvedmath/problems.json}, with \texttt{catalog/problems.json} as the additional-record registry; supplied permanent UnsolvedMath ID \(30001737\), title ``The Converse Unitary Distinction Criterion,'' source problem code \texttt{catalogue entry}. The supplied record belongs to the Oberwolfach Reports Open Problems collection and gives proposed year 2011; no individual proposer is supplied.
\item Supplied report citation: \emph{Automorphic Forms: New Directions} (2011), Oberwolfach Reports, pp.~719--779, DOI 10.4171/owr/2011/14. Supplied source URL and literature-source label ``original OWR report'': \url{https://doi.org/10.4171/owr/2011/14}. The background describes Oberwolfach Reports as documenting mathematical workshops and their open-problem discussions.
\item Supplied literature-source label ``ladder representations'': \url{https://arxiv.org/abs/1411.2420}. No authors, title, or precise applicable theorem are supplied.
\item Supplied literature-source label ``supercuspidal representations'': \url{https://arxiv.org/abs/1909.10450}. No authors, title, or precise applicable theorem are supplied.
\item Supplied literature assessment, checked on \texttt{2026-08-21}: the record labels the problem partially solved and attributes the converse to Jacquet. It reports results for important classes, including generic, ladder, or supercuspidal cases, while describing the general converse as unresolved. This is supplied dated curation derived from the OWR record, not independently verified here; its time-sensitive status and the applicability of the cited results require independent checking.
\item Supplied statement-verification metadata labels the statement \texttt{corrected\_verified} and reports that its mathematical content was checked against the official report, with notation, context, and question form normalized. This records the supplied verification claim; it does not constitute independent source verification in this draft.
\end{itemize}
\end{document}
